\chapter{Arquitectura de servicios}
\label{cap:AnexoB}

Este anexo describe la topología de despliegue del sistema, los servicios que lo componen, sus dependencias y las variantes de configuración disponibles. El sistema se orquesta mediante Docker Compose.

\section{Diagrama de contenedores}

La Figura~\ref{fig:c4-contenedores} muestra la vista de contenedores C4 del sistema: los cuatro componentes desplegables, sus tecnologías y los canales de comunicación entre ellos.

\begin{figure}[ht]
\centering
\resizebox{\textwidth}{!}{% C4 Level 2 — Diagrama de contenedores
% Uso: % C4 Level 2 — Diagrama de contenedores
% Uso: % C4 Level 2 — Diagrama de contenedores
% Uso: \input{figs/c4-contenedores} dentro de figure+\centering
\begin{tikzpicture}[
  every node/.style={font=\small\sffamily},
  actor/.style={
    draw=gray!55, fill=gray!10, thick,
    rectangle, rounded corners=5pt,
    minimum width=2.4cm, minimum height=2.2cm,
    align=center, text width=2.2cm
  },
  cont/.style={
    draw=aquaESI!80, fill=aquaESI!12, thick,
    rectangle, rounded corners=4pt,
    minimum width=3.2cm, minimum height=2.6cm,
    align=center, text width=3.0cm
  },
  db/.style={
    draw=turquesa!80, fill=turquesa!12, thick,
    cylinder, shape border rotate=90,
    minimum height=2.6cm, minimum width=2.8cm,
    align=center, text width=2.4cm
  },
  arr/.style={->, >=Stealth, thick, gray!60},
  biarr/.style={<->, >=Stealth, thick, gray!60},
  lbl/.style={font=\scriptsize\sffamily, fill=white,
              inner sep=1.5pt, draw=gray!30, rounded corners=2pt},
  node distance=1.0cm and 2.2cm
]

\node[actor] (usr) {
  \textbf{Codificador}\\[4pt]
  \small clínico\\[2pt]
  \scriptsize[Actor externo]
};

\node[cont, right=2.0cm of usr] (fe) {
  \textbf{Frontend}\\[6pt]
  \scriptsize Next.js 16, React 19\\TypeScript
};

\node[cont, right=2.2cm of fe] (be) {
  \textbf{Backend}\\[6pt]
  \scriptsize Elixir 1.15, Phoenix 1.8\\Commanded (CQRS/ES)
};

\node[cont, below=1.8cm of be] (ai) {
  \textbf{AI Engine}\\[6pt]
  \scriptsize Python 3.11, FastAPI\\IIC/RigoBERTa-Clinical
};

\node[db, right=2.2cm of be] (pg) {
  \textbf{PostgreSQL 15}\\[6pt]
  \scriptsize Event Store\\+ Proyecciones
};

% Relaciones
\draw[arr] (usr) -- node[lbl, above]{HTTPS} (fe);

\draw[biarr, bend left=20] (fe) to
  node[lbl, above]{\footnotesize REST API} (be);
\draw[biarr, bend right=20] (fe) to
  node[lbl, below]{\footnotesize WebSocket} (be);

\draw[biarr] (be) -- node[lbl, right]{\footnotesize HTTP POST /predict} (ai);

\draw[biarr] (be) -- node[lbl, above]{\footnotesize SQL} (pg);

% Límite del sistema
\begin{pgfonlayer}{background}
  \node[
    draw=aquaESI!60, dashed, thick, rounded corners=10pt,
    fill=aquaESI!4, inner sep=16pt,
    fit=(fe)(be)(ai)(pg),
    label={[aquaESI!80, font=\footnotesize\sffamily\bfseries]above left:%
      Sistema de Codificación CIE-10-ES}
  ] {};
\end{pgfonlayer}

\end{tikzpicture}
 dentro de figure+\centering
\begin{tikzpicture}[
  every node/.style={font=\small\sffamily},
  actor/.style={
    draw=gray!55, fill=gray!10, thick,
    rectangle, rounded corners=5pt,
    minimum width=2.4cm, minimum height=2.2cm,
    align=center, text width=2.2cm
  },
  cont/.style={
    draw=aquaESI!80, fill=aquaESI!12, thick,
    rectangle, rounded corners=4pt,
    minimum width=3.2cm, minimum height=2.6cm,
    align=center, text width=3.0cm
  },
  db/.style={
    draw=turquesa!80, fill=turquesa!12, thick,
    cylinder, shape border rotate=90,
    minimum height=2.6cm, minimum width=2.8cm,
    align=center, text width=2.4cm
  },
  arr/.style={->, >=Stealth, thick, gray!60},
  biarr/.style={<->, >=Stealth, thick, gray!60},
  lbl/.style={font=\scriptsize\sffamily, fill=white,
              inner sep=1.5pt, draw=gray!30, rounded corners=2pt},
  node distance=1.0cm and 2.2cm
]

\node[actor] (usr) {
  \textbf{Codificador}\\[4pt]
  \small clínico\\[2pt]
  \scriptsize[Actor externo]
};

\node[cont, right=2.0cm of usr] (fe) {
  \textbf{Frontend}\\[6pt]
  \scriptsize Next.js 16, React 19\\TypeScript
};

\node[cont, right=2.2cm of fe] (be) {
  \textbf{Backend}\\[6pt]
  \scriptsize Elixir 1.15, Phoenix 1.8\\Commanded (CQRS/ES)
};

\node[cont, below=1.8cm of be] (ai) {
  \textbf{AI Engine}\\[6pt]
  \scriptsize Python 3.11, FastAPI\\IIC/RigoBERTa-Clinical
};

\node[db, right=2.2cm of be] (pg) {
  \textbf{PostgreSQL 15}\\[6pt]
  \scriptsize Event Store\\+ Proyecciones
};

% Relaciones
\draw[arr] (usr) -- node[lbl, above]{HTTPS} (fe);

\draw[biarr, bend left=20] (fe) to
  node[lbl, above]{\footnotesize REST API} (be);
\draw[biarr, bend right=20] (fe) to
  node[lbl, below]{\footnotesize WebSocket} (be);

\draw[biarr] (be) -- node[lbl, right]{\footnotesize HTTP POST /predict} (ai);

\draw[biarr] (be) -- node[lbl, above]{\footnotesize SQL} (pg);

% Límite del sistema
\begin{pgfonlayer}{background}
  \node[
    draw=aquaESI!60, dashed, thick, rounded corners=10pt,
    fill=aquaESI!4, inner sep=16pt,
    fit=(fe)(be)(ai)(pg),
    label={[aquaESI!80, font=\footnotesize\sffamily\bfseries]above left:%
      Sistema de Codificación CIE-10-ES}
  ] {};
\end{pgfonlayer}

\end{tikzpicture}
 dentro de figure+\centering
\begin{tikzpicture}[
  every node/.style={font=\small\sffamily},
  actor/.style={
    draw=gray!55, fill=gray!10, thick,
    rectangle, rounded corners=5pt,
    minimum width=2.4cm, minimum height=2.2cm,
    align=center, text width=2.2cm
  },
  cont/.style={
    draw=aquaESI!80, fill=aquaESI!12, thick,
    rectangle, rounded corners=4pt,
    minimum width=3.2cm, minimum height=2.6cm,
    align=center, text width=3.0cm
  },
  db/.style={
    draw=turquesa!80, fill=turquesa!12, thick,
    cylinder, shape border rotate=90,
    minimum height=2.6cm, minimum width=2.8cm,
    align=center, text width=2.4cm
  },
  arr/.style={->, >=Stealth, thick, gray!60},
  biarr/.style={<->, >=Stealth, thick, gray!60},
  lbl/.style={font=\scriptsize\sffamily, fill=white,
              inner sep=1.5pt, draw=gray!30, rounded corners=2pt},
  node distance=1.0cm and 2.2cm
]

\node[actor] (usr) {
  \textbf{Codificador}\\[4pt]
  \small clínico\\[2pt]
  \scriptsize[Actor externo]
};

\node[cont, right=2.0cm of usr] (fe) {
  \textbf{Frontend}\\[6pt]
  \scriptsize Next.js 16, React 19\\TypeScript
};

\node[cont, right=2.2cm of fe] (be) {
  \textbf{Backend}\\[6pt]
  \scriptsize Elixir 1.15, Phoenix 1.8\\Commanded (CQRS/ES)
};

\node[cont, below=1.8cm of be] (ai) {
  \textbf{AI Engine}\\[6pt]
  \scriptsize Python 3.11, FastAPI\\IIC/RigoBERTa-Clinical
};

\node[db, right=2.2cm of be] (pg) {
  \textbf{PostgreSQL 15}\\[6pt]
  \scriptsize Event Store\\+ Proyecciones
};

% Relaciones
\draw[arr] (usr) -- node[lbl, above]{HTTPS} (fe);

\draw[biarr, bend left=20] (fe) to
  node[lbl, above]{\footnotesize REST API} (be);
\draw[biarr, bend right=20] (fe) to
  node[lbl, below]{\footnotesize WebSocket} (be);

\draw[biarr] (be) -- node[lbl, right]{\footnotesize HTTP POST /predict} (ai);

\draw[biarr] (be) -- node[lbl, above]{\footnotesize SQL} (pg);

% Límite del sistema
\begin{pgfonlayer}{background}
  \node[
    draw=aquaESI!60, dashed, thick, rounded corners=10pt,
    fill=aquaESI!4, inner sep=16pt,
    fit=(fe)(be)(ai)(pg),
    label={[aquaESI!80, font=\footnotesize\sffamily\bfseries]above left:%
      Sistema de Codificación CIE-10-ES}
  ] {};
\end{pgfonlayer}

\end{tikzpicture}
}
\caption{Diagrama de contenedores C4 del sistema CIE-10-ES.}
\label{fig:c4-contenedores}
\end{figure}

\section{Topología de contenedores}

El \emph{Compose} base define cinco servicios Docker~\cite{merkel2014docker} y construye la imagen base de ML (\texttt{ml-base}), que solo se construye y no se ejecuta. La Tabla~\ref{tab:servicios} los resume, junto con \texttt{mailpit} (definido en \texttt{docker-compose.mailpit.yml} y solo de desarrollo), con su función, tecnología, puerto expuesto y dependencias. Los ficheros de proxy y de monitorización añaden otros diez servicios de apoyo (Tabla~\ref{tab:servicios-apoyo}), que no forman parte del sistema de codificación y por eso no aparecen en la Figura~\ref{fig:c4-contenedores}; su papel en el despliegue se describe en el Anexo~\ref{cap:AnexoI} y la monitorización, en el Anexo~\ref{cap:AnexoK}.

\begin{table}[h]
\centering
\caption{Servicios del sistema y sus características de despliegue.}
\label{tab:servicios}
\begin{tabular}{lP{3.2cm}rP{2.3cm}P{2.8cm}}
\hline
\textbf{Servicio} & \textbf{Tecnología} & \textbf{Puerto} & \textbf{Depende de} & \textbf{Rol} \\
\hline
\texttt{frontend} & Next.js 16, React 19, TypeScript & 3000 & backend & Interfaz web \\
\texttt{backend} & Elixir 1.19, Phoenix 1.8 & 4000 & db, ai\_engine & Orquestador \\
\texttt{ai\_engine} & Python 3.11, FastAPI & 8000 & db & Inferencia IA \\
\texttt{db} & PostgreSQL 18 & 5432 &  ninguna & Persistencia \\
\texttt{mailpit} & Mailpit & 8025/1025 &  ninguna & Correo de captura (solo desarrollo) \\
\texttt{ml-base} & CUDA + PyTorch + Transformers & ninguno &  ninguna & Imagen base de ML (solo se construye) \\
\texttt{training} & Python, Jupyter Lab & 8888 & ml-base & Entrenamiento \\
\hline
\end{tabular}
\end{table}

\begin{table}[h]
\centering
\caption{Servicios de apoyo (proxy y monitorización). Fuente: \texttt{docker-compose-proxy.yml} y \texttt{docker-compose.monitoring.yml}.}
\label{tab:servicios-apoyo}
\begin{tabular}{llP{7.2cm}}
\hline
\textbf{Servicio} & \textbf{Fichero} & \textbf{Función} \\
\hline
\texttt{traefik} & proxy & Proxy inverso: enruta por nombre de host a frontend, backend, Grafana y a su propio panel \\
\texttt{cloudflared} & proxy & Túnel de Cloudflare: entrada del tráfico sin abrir puertos \\
\texttt{maintenance} & proxy & Página de mantenimiento (nginx) cuando el frontend no responde \\
\texttt{prometheus} & monitorización & Recoge y almacena métricas \\
\texttt{grafana} & monitorización & Paneles de visualización \\
\texttt{loki} & monitorización & Almacén de registros \\
\texttt{promtail} & monitorización & Envía los registros de los contenedores a Loki \\
\texttt{cadvisor} & monitorización & Métricas de uso de los contenedores \\
\texttt{blackbox\_exporter} & monitorización & Sondas de disponibilidad de los servicios \\
\texttt{node\_exporter} & monitorización & Métricas del equipo anfitrión \\
\hline
\end{tabular}
\end{table}

\section{Flujo de arranque}

El orden de arranque impuesto por las dependencias de Docker Compose es el siguiente:

\begin{enumerate}
 \item \textbf{PostgreSQL} (\texttt{db}) arranca primero; expone el puerto 5432 internamente.
 \item \textbf{AI Engine} arranca cuando la base de datos está lista (\texttt{service\_healthy}); carga el modelo en memoria durante el arranque (\textasciitilde{}13\,GB de RAM del sistema en modo CPU).
 \item \textbf{Backend} espera a que \texttt{db} y \texttt{ai\_engine} estén disponibles; ejecuta las migraciones Ecto al arrancar y crea las tablas del \emph{event store} si no existen.
 \item \textbf{Frontend} arranca una vez el backend está activo.
 \item \textbf{Mailpit} arranca en paralelo; solo se usa en entornos de desarrollo.
\end{enumerate}

\section{Variantes de despliegue}

El repositorio incluye tres ficheros de configuración adicional de Docker Compose que adaptan el despliegue al entorno de ejecución:

\begin{description}
 \item[\texttt{docker-compose.gpu.yml}] Configuración por defecto para máquinas con GPU NVIDIA. El servicio \texttt{ai\_engine} arranca con \texttt{DEVICE=cuda} y acceso al \emph{runtime} NVIDIA. Requiere el controlador NVIDIA y el toolkit de contenedores de NVIDIA instalados en el anfitrión. El modelo se carga en VRAM (verificado con una RTX~4080 de 16\,GB), lo que libera la mayor parte de los \textasciitilde{}13\,GB de RAM del sistema que ocuparía en modo CPU. La construcción de la imagen GPU incluye la compilación de flash attention, lo que requiere al menos 32\,GB de RAM durante el proceso de compilación. Tiempo de inferencia: menos de 1 segundo por informe.

 \item[\texttt{docker-compose.cpu.yml}] Configuración para máquinas sin GPU. Sobreescribe \texttt{DEVICE=cpu} en \texttt{ai\_engine}. El backend de atención cambia automáticamente a SDPA (\emph{Scaled Dot-Product Attention}). Requiere al menos 16\,GB de RAM del sistema; el entorno completo con todos los modelos cargados ocupa aproximadamente 13\,GB. Tiempo de inferencia: menos de 1 segundo por informe en hardware moderno (verificado en Ryzen 7900X y Apple M4).

 \item[\texttt{docker-compose.mock.yml}] Sustituye el servicio \texttt{ai\_engine} por \texttt{ai\_engine\_mock}, una implementación ligera en FastAPI que devuelve predicciones ficticias sin cargar ningún modelo. Recomendado para el desarrollo del frontend y del backend cuando no se necesita inferencia real.
\end{description}

\section{Variables de entorno}

La Tabla~\ref{tab:envvars} recoge las variables de entorno más relevantes del sistema, con sus valores por defecto en el entorno de desarrollo. El fichero \texttt{.env.example} del repositorio contiene el listado completo.

\begin{table}[!t]
\centering
\caption{Variables de entorno del sistema.}
\label{tab:envvars}
\begin{tabular}{lP{5.2cm}l}
\hline
\textbf{Variable} & \textbf{Valor por defecto} & \textbf{Consumidor} \\
\hline
\multicolumn{3}{l}{\textit{Base de datos}} \\
\texttt{POSTGRES\_USER} & \texttt{postgres} & db \\
\texttt{POSTGRES\_PASSWORD} & \texttt{changeme} & db, backend \\
\texttt{DATABASE\_URL} & \nolinkurl{postgresql://postgres:changeme@db:5432/cie10_app} & backend \\
\texttt{EVENT\_STORE\_URL} & \nolinkurl{postgresql://postgres:changeme@db:5432/cie10_eventstore} & backend \\
\hline
\multicolumn{3}{l}{\textit{Backend}} \\
\texttt{SECRET\_KEY\_BASE} & (mínimo 64 bytes en prod.) & backend \\
\texttt{MIX\_ENV} & \texttt{dev} & backend \\
\texttt{BACKEND\_PORT} & \texttt{4000} & backend \\
\hline
\multicolumn{3}{l}{\textit{Frontend}} \\
\texttt{NEXT\_PUBLIC\_API\_URL} & \texttt{http://\allowbreak{}localhost:4000/api} & frontend \\
\texttt{NEXT\_PUBLIC\_WS\_URL} & \texttt{ws://\allowbreak{}localhost:4000/socket} & frontend \\
\texttt{NODE\_ENV} & \texttt{development} & frontend \\
\hline
\multicolumn{3}{l}{\textit{Motor de IA}} \\
\texttt{AI\_ENGINE\_PORT} & \texttt{8000} & ai\_engine \\
\texttt{AI\_ENGINE\_URL} & \texttt{http://\allowbreak{}ai\_engine:8000} & backend \\
\texttt{SUMMARIZER\_MODEL} & \texttt{none} & ai\_engine \\
\texttt{SUMMARIZER\_MODE} & \texttt{paraphrase} & ai\_engine \\
\texttt{HF\_TOKEN} & (token HuggingFace) & ai\_engine, training \\
\hline
\multicolumn{3}{l}{\textit{Producción y red}} \\
\texttt{DOMAIN} & \texttt{clicoder.app} & traefik \\
\texttt{CLOUDFLARE\_TUNNEL\_TOKEN} & (token del túnel) & cloudflared \\
\texttt{CORS\_ORIGINS} & \texttt{http://\allowbreak{}localhost:3000,\allowbreak{}http://\allowbreak{}localhost:5173} & backend \\
\hline
\multicolumn{3}{l}{\textit{Observabilidad}} \\
\texttt{GRAFANA\_PASSWORD} & \texttt{changeme} & grafana \\
\texttt{SENTRY\_DSN} & (DSN de Sentry) & backend \\
\texttt{SENTRY\_DSN\_AI} & (DSN de Sentry) & ai\_engine \\
\hline
\multicolumn{3}{l}{\textit{Datos iniciales}} \\
\texttt{SEED\_ADMIN\_EMAIL} & \texttt{myadminuser@clicoder.app} & backend \\
\texttt{SEED\_ADMIN\_PASSWORD} & \texttt{password123!} & backend \\
\hline
\multicolumn{3}{l}{\textit{Email}} \\
\texttt{SMTP\_HOST} & \texttt{mailpit} (desarrollo; en despliegue, el servidor SMTP de Mailjet) & backend \\
\texttt{SMTP\_PORT} & \texttt{1025} (desarrollo) & backend \\
\texttt{APP\_HOST} & \texttt{localhost} & backend \\
\hline
\end{tabular}
\end{table}

\section{Estructura del repositorio}

\begin{description}
 \item[\texttt{ai\_engine/}] Servicio FastAPI de inferencia. Contiene \texttt{main.py} (punto de entrada), \texttt{classifier.py} (clase \texttt{CIE10Classifier}), \texttt{train.py} (entrenamiento), \texttt{baseline\_dict.py} (\emph{fallback} por diccionario) y \texttt{model/} (artefactos del modelo entrenado).
 \item[\texttt{ai\_engine\_mock/}] Implementación simulada del motor de IA para desarrollo.
 \item[\texttt{backend/}] Aplicación Phoenix. \texttt{lib/app/} contiene la lógica de negocio (agregados, comandos, eventos, proyectores); \texttt{lib/app\_web/} contiene la capa HTTP (controladores, canales, \emph{plugs}, router).
 \item[\texttt{frontend/}] Aplicación Next.js. \texttt{src/components/} contiene los componentes de UI; \texttt{src/contexts/} los contextos de estado; \texttt{src/lib/} la configuración y utilidades de red.
 \item[\texttt{training/}] Entorno de entrenamiento aislado. \texttt{csv\_import\_scripts/} para descarga de datos; \texttt{data\_analysis/} para análisis exploratorio; \texttt{bert-classifier/} para notebooks de entrenamiento.
 \item[\texttt{docker/grafana/}] Configuración de aprovisionamiento de Grafana: paneles de monitorización del sistema (métricas de contenedores, backend y motor de IA) listos para importar automáticamente al arrancar el stack de monitorización.
 \item[\texttt{docker/prometheus/}] Configuración de Prometheus (\texttt{prometheus.yml}): definición de los \emph{targets} de scraping para el backend, el motor de IA, cAdvisor y Node Exporter.
 \item[\texttt{tfg/}] Contiene los archivos del trabajo de fin de grado, incluyendo capítulos, anexos y figuras.
\end{description}
